\documentclass[10pt]{article}
\usepackage[a4paper,margin=1.8cm]{geometry}
\usepackage{amsmath,amssymb,booktabs,enumitem}
\usepackage[hidelinks]{hyperref}
\setlist{nosep,leftmargin=1.3em}
\begin{document}
\begin{center}
{\Large\bfseries Climbing to the Frontier --- METHODOLOGY}\\[2pt]
BAINSA Hackathon, Track 3 (Anthropic), 26 September 2026 --- Team: Lorenzo Pazienza (lead), Seby, Teo, Niccol\`o
\end{center}

\section*{1. Game plan}
\begin{itemize}
\item \textbf{Triage first.} Before writing anything we read all 24 cells, checked the literature for each problem (Griggs--Ho for P3, O'Bryant and Fornal--Sun for P4, Bilyk--Matzke and Fodor--V\'igh--Zarn\'ocz for P1), and ranked the cells by points against feasibility. The cheap cells were closed first. The time-boxed effort went to the C5 cells, which carry the weight, and partial C6 work came last.
\item \textbf{One problem per person, one agent per sub-task.} Each team member owned a problem and ran Claude Code (Opus) sessions for it. The lead ran several further sessions (P2 search, P4 cross-check, P4~C6, and a finaliser that assembles the hand-in), Codex as an independent second model, and a Claude coordinator that tracked the state of every cell. A single status table (cell $\to$ file $\to$ state $\to$ owner) was the source of truth, and the lead kept the plan and made every accept/reject decision.
\item \textbf{Computation only where the rules allow it.} We used the computer to explore, and relied on it in a proof only over a finite set that our own argument had reduced to, with code shipped and runtimes stated. Every computational cell follows the five-point exhaustiveness certificate.
\end{itemize}

\section*{2. Correctness: nothing is accepted on one model's word}
\begin{itemize}
\item \textbf{Exact evaluators outside the model.} Scores such as the uphill-path counts come from a small evaluator with two independent counting methods, a DP and brute-force enumeration. No number in a write-up comes from the model's own arithmetic.
\item \textbf{Two independent implementations for every certificate.} In P4 the SAT pipeline (Codex) and the DFS pipeline (Teo) were written separately. Their survivor lists agree element by element, with the same SHA-256, for all $k\le16$. In P2 the SAT lemmas were run with CaDiCaL and Glucose, and re-implemented from scratch by a different agent that shared no code with the first. Sanity runs at the feasible sizes show the encodings are not vacuous.
\item \textbf{Adversarial review by a different agent than the author.} Examples:
  \begin{itemize}
  \item P1~C5 (the cycle lemma) was reviewed line by line by Codex and by Claude, with numerical stress tests (40k random cycles, $10^6$ points for the key inequality).
  \item P3~C4/C5 were re-verified by an independent exhaustive checker that finds cycles straight from the definition, without assuming Brandt's theorem, up to $k=12$, beyond the authors' $n\le60$.
  \item P2's new lower-bound proof was checked step by step, including recomputing the Delsarte table.
  \item A final read-only referee pass (Codex) covered all cells before hand-in. It found packaging and labelling problems (labelling-file format, code not attached to the upload, a ``by hand'' label on a computer-assisted step, missing certificate data for P4), and we fixed them before submitting.
  \end{itemize}
\item \textbf{Honest labels.} A claim is \emph{solved} only with a complete proof. Everything else is \emph{partial}, with the proved part stated. Every result that depends on a computer says \emph{computer-assisted}, and unfinished runs are reported as unfinished. For example, we claim the certified boundary $K^\ast=15$ rather than 16, because the second implementation had not finished deciding $k=16$. We also say explicitly that a planned $k=14$ rerun was not done.
\end{itemize}

\section*{3. Efficiency (compute and tokens)}
\begin{itemize}
\item \textbf{Measure before waiting.} We estimated the runtime of long jobs from uniform random samples of their inputs. That showed a $k=44$ certificate needs about 240 CPU-hours and $k=48$ even more, so we stopped both and kept the (computer-assisted, 4-second) proof of $k=32,38,42$.
\item \textbf{Stop duplicates early.} When two sessions were found working on the same cell (P1~C3; the $k=32$ case of P4~C6), one was stopped and redirected.
\item \textbf{Short, targeted prompts to agents.} Agents received the official cell text, the exact files to read and a deadline, and wrote into their own folders. Only the finaliser wrote into the hand-in folder. This avoided re-reading large contexts and conflicting edits.
\end{itemize}

\section*{4. Human verification and intervention}
\begin{itemize}
\item The lead read every agent report and decided what to accept, what to send for review, and what to stop. Several agent claims were corrected before hand-in, for instance an over-claimed status on P2 and a mislabelled ``solved'' on P3~C3/C5, which became \emph{partial}.
\item Where an agent had paraphrased a cell statement, we replaced it with the official text from the problem page (for example P4~C4/C5), and we relabelled cells whose write-up claimed more than it proved.
\item When a teammate's independent work turned out stronger than ours (the P2 lower bound via decycling sets), we switched the hand-in to it after verifying it, rather than keeping our own.
\end{itemize}

\section*{5. What went wrong, and the fix}
Around midday several agents and people were working on the same cells, from four copies of the same folder. We fixed this with three rules: one owner per cell, one hand-in folder written by one session only, and a live status table. From then on no cell had two writers.
\end{document}
